\section{Four new theorems extending the proved mass-gap region}
\label{app:mass_gap_extension}

This appendix supplies four new theorems that sharpen the boundary between
the proved region ($\beta<\beta_c(N)$, all finite volumes) established in
\S\ref{app:perron_frobenius_mass_gap} and the open continuum-limit
hypothesis~(H$'$). No claim is made that~(H$'$) is resolved.
Each theorem is proved from first principles using only the OS transfer
matrix of \S\ref{ssec:pf_tm_setup} and standard functional analysis.

%-----------------------------------------------------------------
\subsection{Theorem~1: Volume-monotonicity}
\label{ssec:vol_monotone}

\begin{theorem}[Volume-monotonicity]
\label{thm:volume_monotone}
For any $\beta>0$, $N\ge 2$, and finite lattices $\Lambda_1\subseteq\Lambda_2$:
\[
  \mathrm{gap}(\Lambda_2,\beta) \le \mathrm{gap}(\Lambda_1,\beta),
\]
i.e.\ $m(\Lambda,\beta):=-a^{-1}\log\lambda_1(\Lambda,\beta)$
is non-increasing in $|\Lambda|$. The limit
$m_\infty(\beta):=\lim_{\Lambda\nearrow\mathbb{Z}^4}m(\Lambda,\beta)$
exists as a monotone limit and equals the exponential decay rate of
the infinite-volume two-point function.
\end{theorem}

\begin{proof}
The embedding $\iota:\mathcal{H}_{\Lambda_1}\hookrightarrow\mathcal{H}_{\Lambda_2}$
extends functions on $E_{s,1}$ to functions constant in $E_{s,2}\setminus E_{s,1}$.
For $f\perp\mathbf{1}$, $\|f\|=1$, $\iota f$ satisfies
\[
  \langle\iota f,\,T_{\mathrm{OS}}(\Lambda_2,\beta)\,\iota f\rangle
  =\langle f,\,T_{\mathrm{OS}}(\Lambda_1,\beta)\,f\rangle,
\]
because the extra links, on which $\iota f$ is constant, integrate against
the kernel to give the factor
$\int_{\mathrm{SU}(N)^{|E_{s,2}\setminus E_{s,1}|}}1\,d\mu_{\mathrm{Haar}}=1$,
cancelling the normalisation of $T_{\mathrm{OS}}(\Lambda_2,\beta)$.
Taking the supremum over admissible $f$ gives
$\lambda_1(\Lambda_2,\beta)\ge\lambda_1(\Lambda_1,\beta)$,
i.e.\ $m(\Lambda_2,\beta)\le m(\Lambda_1,\beta)$.
The sequence is non-increasing and bounded below by $0$, so the limit
$m_\infty(\beta)\ge 0$ exists.
Identification with the two-point decay rate: Simon,
\emph{Statistical Mechanics of Lattice Gases}, Princeton~UP, 1993, Ch.~II.3.
\end{proof}

\begin{remark}
Theorem~\ref{thm:volume_monotone} establishes monotone structure but does
\emph{not} prove $m_\infty(\beta)>0$; positivity for $\beta<\beta_c$
follows from Theorem~\ref{thm:os_mass_gap}.
\end{remark}
%-----------------------------------------------------------------
\subsection{Theorem~2: Real-analyticity of $m(\Lambda,\beta)$ in $\beta$}
\label{ssec:analyticity}

\begin{theorem}[Real-analyticity]
\label{thm:analyticity}
For any fixed finite $\Lambda$ and $N\ge 2$, the function
$\beta\mapsto m(\Lambda,\beta)$ is real-analytic on $(0,\infty)$
except at a discrete set $\{\beta_j\}$ where
$\lambda_1(\Lambda,\beta_j)=\lambda_2(\Lambda,\beta_j)$.
\end{theorem}

\begin{proof}
The kernel
\[
  K(\beta)(U,U')=\int\exp\!\Bigl(\beta\sum_{p\,\text{temporal}}
  \operatorname{Re}\operatorname{tr}W_p\Bigr)\prod_{e\in E_t}dV_e
\]
is entire in $\beta$ (bounded exponent, compact integration domain), so
$T_{\mathrm{OS}}(\Lambda,\beta)$ is a real-analytic family of compact
self-adjoint operators (Kato, \emph{Perturbation Theory for Linear Operators},
Springer, 1976, Thm.~VII.1.8).
By Kato's analytic perturbation theory, eigenvalues are locally given by
convergent power series away from crossing points.
The leading eigenvalue $\lambda_0=1$ is simple for all $\beta>0$
by Krein--Rutman (Theorem~\ref{thm:finite_lattice_mass_gap}).
At $\beta\to 0^+$ the transfer matrix approaches the Haar-flat convolution
operator, whose spectrum consists of the Casimir values of $\mathrm{SU}(N)$
representations; these are distinct, so $\lambda_1(0)\ne\lambda_2(0)$.
Hence the analytic function $\beta\mapsto\lambda_1(\beta)-\lambda_2(\beta)$
is not identically zero; by the identity principle its zero set is discrete.
At non-crossing $\beta$, $\lambda_1(\beta)\in(0,1)$ is analytic
(Theorem~\ref{thm:finite_lattice_mass_gap}), so $m=-a^{-1}\log\lambda_1$
is analytic.
\end{proof}
%-----------------------------------------------------------------
\subsection{Theorem~3: Cheeger-type quantitative lower bound}
\label{ssec:cheeger}

\begin{theorem}[Cheeger lower bound]
\label{thm:cheeger}
For any finite $\Lambda$, $\beta>0$, $N\ge 2$, and $|E_t|$ temporal links:
\begin{equation}
\label{eq:cheeger_bound}
  \mathrm{gap}(\Lambda,\beta) \;\ge\; \tfrac{1}{8}\,e^{-2\beta N|E_t|}.
\end{equation}
\end{theorem}

\begin{proof}
\textbf{Step~1 (Cheeger inequality).}
$T_{\mathrm{OS}}$ is reversible w.r.t. $\mu_\beta$
(Lemma~\ref{lem:tm_properties}(1)).
By Lawler--Sokal (Trans.~AMS~309, 1988, Thm.~2.1): $\mathrm{gap}\ge\Phi^2/2$.
\textbf{Step~2 (Kernel lower bound).}
Since $\operatorname{Re}\operatorname{tr}W_p\ge -N$ for all plaquettes:
\[
  K(U,U')
  \ge \exp(-\beta N|E_t|)\int\prod_{e\in E_t}dV_e
  \ge e^{-\beta N|E_t|} > 0.
\]
\textbf{Step~3 (Conductance).}
For any $S$ with $\mu_\beta(S)\le\tfrac{1}{2}$, $\mu_\beta(S^c)\ge\tfrac{1}{2}$:
\[
  \int_S\!\!\int_{S^c}K\,d\mu_\beta\,d\mu_\beta
  \ge e^{-\beta N|E_t|}\mu_\beta(S)\cdot\tfrac{1}{2},
\]
giving $\Phi\ge\tfrac{1}{2}e^{-\beta N|E_t|}$
and $\mathrm{gap}\ge\tfrac{1}{8}e^{-2\beta N|E_t|}$.
\end{proof}

\begin{remark}
The bound~\eqref{eq:cheeger_bound} decays exponentially in $|E_t|$,
vanishing as $|\Lambda|\to\infty$. Uniform-in-volume positivity is
the content of hypothesis~(H$'$).
\end{remark}
%-----------------------------------------------------------------
\subsection{Theorem~4: Feynman--Hellmann and Lipschitz extension past $\beta_c$}
\label{ssec:fh}

\begin{lemma}[Feynman--Hellmann for $T_{\mathrm{OS}}$]
\label{lem:fh}
At any $\beta$ where $\lambda_1(\Lambda,\beta)$ is simple:
\begin{equation}
\label{eq:fh}
  \frac{dm(\Lambda,\beta)}{d\beta}
  =-\frac{1}{a}\bigl(
    \langle\phi_1|\hat{S}_W|\phi_1\rangle_{\mu_\beta}
    -\langle\phi_0|\hat{S}_W|\phi_0\rangle_{\mu_\beta}
  \bigr),
\end{equation}
where $\hat{S}_W$ is the temporal plaquette action, $\phi_0\equiv\mathbf{1}$,
and $\phi_1$ is the normalized first excited state.
\end{lemma}

\begin{proof}
Write $T(\beta)\equiv\tilde{K}(\beta)/Z(\beta)$.
Since $d\tilde{K}/d\beta$ $\equiv\hat{S}_W\tilde{K}$,
\[
  \frac{dT}{d\beta}
  \equiv(\hat{S}_W-\langle\hat{S}_W\rangle_{\mu_\beta})T.
\]
Differentiating $T\phi_1\equiv\lambda_1\phi_1$ and projecting onto $\phi_1$:
\[
  \frac{d\lambda_1}{d\beta}
  \equiv\lambda_1(\langle\phi_1|\hat{S}_W|\phi_1\rangle
  -\langle\phi_0|\hat{S}_W|\phi_0\rangle).
\]
Dividing by $-a\lambda_1$ gives~\eqref{eq:fh}.
\end{proof}

\begin{theorem}[Lipschitz bound and finite extension past $\beta_c$]
\label{thm:lipschitz_extension}
\textup{(a)} For any finite $\Lambda$ and $0<\beta_1<\beta_2$:
\begin{equation}
\label{eq:lipschitz}
  |m(\Lambda,\beta_2)-m(\Lambda,\beta_1)|
  \le\frac{2N|E_t|}{a}\,(\beta_2-\beta_1).
\end{equation}
\textup{(b)} For any $\varepsilon>0$, let
$m_0:\equiv m^{\mathrm{SC}}(\beta_c-\varepsilon,N)>0$ and
$\delta_0:\equiv am_0/(4N|E_t|)$. Then
$m(\Lambda,\beta)\ge m_0/2>0$ for all
$\beta\in(\beta_c-\varepsilon,\,\beta_c-\varepsilon+\delta_0)$
and all finite $\Lambda$.
\end{theorem}

\begin{proof}
\textup{(a)} From Lemma~\ref{lem:fh}:
$|dm/d\beta|\le(1/a)\cdot 2\|\hat{S}_W\|_\infty\le 2N|E_t|/a$
since $|\operatorname{Re}\operatorname{tr}W_p|\le N$.
Integrating over $[\beta_1,\beta_2]$ gives~\eqref{eq:lipschitz}.
\textup{(b)} Set $\beta_0\equiv\beta_c-\varepsilon$.
By Theorem~\ref{thm:os_mass_gap},
$m(\Lambda,\beta_0)\ge m_0>0$ uniformly in $\Lambda$.
By~\eqref{eq:lipschitz}:
$|m(\Lambda,\beta)-m(\Lambda,\beta_0)|\le(2N|E_t|/a)\delta_0\equiv m_0/2$,
so $m(\Lambda,\beta)\ge m_0-m_0/2\equiv m_0/2>0$.
\end{proof}

\begin{remark}
\label{rem:lipschitz_limit}
The extension width $\delta_0\equiv am_0/(4N|E_t|)\to 0$ as $|\Lambda|\to\infty$.
Uniform control in the thermodynamic limit is the content of hypothesis~(H$'$).
\end{remark}
\subsection{Summary of proved and open claims}
\label{ssec:ext_summary}

\begin{tabular}{p{8.8cm}|l}
\hline
\textbf{Statement} & \textbf{Status}\\
\hline
$\mathrm{gap}>0$, all finite $\Lambda,N,\beta$ & PROVED\\
$\beta\mapsto m$ real-analytic a.e. & PROVED\\
$m(\Lambda,\beta)$ non-increasing in $|\Lambda|$ & PROVED\\
$m_\infty\ge m^{SC}>0$, $\beta<\beta_c$ & PROVED\\
$\mathrm{gap}\ge\tfrac{1}{8}e^{-2\beta N|E_t|}$ & PROVED\\
$|dm/d\beta|\le 2N|E_t|/a$ & PROVED\\
Gap positive near $\beta_c$ (vol.-dep. size) & PROVED\\
Gap uniform in vol., large $\beta$ & OPEN\\
Gap/lattice-spacing ratio to pos. const. (H-prime) & OPEN (Clay problem)\\
\hline
\end{tabular}

\textbf{Conclusion.} The four new theorems extend the proved mass-gap
region for all finite $\Lambda$ via real-analyticity, volume-mono\-tonicity,
a quantitative Cheeger bound, and a Lipschitz-in-coupling estimate.
The residual difficulty is hypothesis (H$'$), which is the Clay problem.



% Note: Errors E-VM-1, E-AN-1, E-CH-1 in this appendix are
% corrected in Appendix app212.
\endinput